Our recent efforts to improve the reliability of spoofing countermeasures for automatic speaker verification (ASV) have confirmed that: (i)~artefacts from converted voice or synthetic speech spoofing attacks reside at the sub-band level~\cite{odyssey2020CQCC}; (ii)~detection performance can be improved through the use of high-spectral-resolution front-ends, particularly those which emphasize information within the same sub-bands~\cite{IS2020LFCC}. When used even with a trivial Gaussian mixture model (GMM) back-end, high-spectral-resolution front-ends lead to competitive performance that would have obtained fourth position (out of 48 submissions) in the ASVspoof 2019 challenge. 

Unfortunately, though, some spoofing attacks continue to evade detection, particularly the infamous A17 attack, a neural network voice conversion (VC) based attack that uses a generalised direct waveform modification method for waveform generation. 
%Using this method, a F0 transformed residual signal is generated by applying an overlap-add technique based on waveform similarity (WSOLA), inverse filtering and resampling to the input speech and then converted speech is generated by filtering the residual signal with a synthesis filter using the mel-log spectrum approximation (MLSA) filter) within the ASVspoof 2019 database.
While we have seen that even A17 attacks can be detected when classifiers are learned with representative data (i.e.\ through learning with evaluation data), performance is poor when the same classifier is learned using only standard training data. Hence, while high-spectral-resolution front-ends \emph{do} capture A17 artefacts, they \emph{cannot} be detected.

The same observation shows that, once training data for a newly discovered spoofing attack is available, then vulnerabilities can be fixed easily simply by re-learning with representative training data. 
%This is akin to waiting for vulnerabilities to be exposed before taking retrospective action. 
More proactive or preemptive strategies that offer some protection against the unexpected are preferable. This idea embodies the objectives of the ASVspoof initiative since its inception, namely the pursuit of generalisable spoofing countermeasures that are capable of detecting attacks not seen in training data. The latter are characteristic of in-the-wild settings in which the nature of spoofing attacks cannot be predicted and will evolve continuously. The question addressed in this paper is whether A17 attacks can be detected even in the \emph{absence} of representative training data.

This is a difficult question to answer reliably in a post-evaluation setting where the characteristics of attack A17 are now known and cannot be forgotten. One-class classifiers trained only on bona fide data without the use of any spoofed data are an obvious candidate solution and have shown some previous success in anomaly, novelty and spoofing detection tasks~\cite{EURECOM+4093,fatemifar2019spoofing,fatemifar2019combining,ohki2019efficient}. Anecdotal evidence, however, shows that many, including the authors of the current article, have failed to apply them successfully to ASVspoof 2019 data. Since A17 attacks cannot be detected using our current countermeasure (at least not without representative A17 training data) we must seek cues elsewhere. 

Our assumption is that the use of hand-crafted features does not offer the best potential to detect unforeseen attacks since they rely too much upon the characterisation of artefacts corresponding to known attacks. Attack-specific artefacts are likely to be insufficient for the detection of unseen attacks and a higher-level, more generalisable representation is needed to ensure robust performance. This paper reports what is, to the best of our knowledge, the first application of RawNet2~\cite{jung2020improved} deep neural network classifiers to anti-spoofing. RawNet2 operates directly upon the raw speech waveform, circumventing almost entirely any dependence upon hand-crafted features. The goal of this work is not just to design more reliable spoofing countermeasures, but also to determine whether such end-to-end approaches that utilise some degree of automatic feature learning can improve performance in a \emph{worst-case scenario} such as that of A17.

The paper is organised as follows. The original RawNet2 architecture and its application to ASV are described in Section~2. Modifications made to RawNet2 in order to apply it successfully to anti-spoofing are described in Section~3. Experimental work is described in Section~4. Our conclusions are presented in Section~5.